\begin{table}[t]
\centering
\caption{A/B/C observed outcomes and known completed-response resources. C has two unavailable outcomes; its primary observed rate is 1/58. Repair means use the 60 eligible-case denominator, including zero completed repairs. Resource totals include the shared initial in each counterfactual row.}
\label{tab:abc}
\small
\begin{tabular}{lrrrr}
\toprule
Condition & Verified/observed & Missing & Rate & Wilson 95\% \\
\midrule
A & 0/60 & 0 & 0.00\% & [0.00, 6.02] \\
B & 1/60 & 0 & 1.67\% & [0.29, 8.86] \\
C & 1/58 & 2 & 1.72\% & [0.31, 9.14] \\
\bottomrule
\end{tabular}
\par\vspace{4pt}
\begin{tabular}{lrrrr}
\toprule
Cond. & Repairs (mean) & Tokens in/out & Call wall s & Lean s \\
\midrule
A & 0 (0.00) & 12,929/27,862 & 1908.6 & 2361.9 \\
B & 178 (2.97) & 122,029/83,485 & 8120.5 & 8923.9 \\
C & 172 (2.87) & 252,935/74,778 & 9870.8 & 8906.3 \\
\bottomrule
\end{tabular}
\par\vspace{4pt}
\begin{minipage}{.98\linewidth}\footnotesize
Initial success is 0/60. On the 58 complete B/C pairs each method succeeds once: paired difference 0.00 percentage points. B consumes 178 logical repair intents; C consumes 174, of which 172 complete. Sensitivity counting missing C as failure is 1/60 (1.67\%), not the observed C rate. Call wall times exclude the unavailable response and preflight rejection; they are not full-run cost estimates.
\end{minipage}
\end{table}
